\section{Partial envelopes on a weighted compact set}
\label{sec:hermitian-partial-envelopes}

A continuous obstacle on a non-pluripolar compact set bounds all admissible potentials globally. Imposing a multiplier-ideal constraint then fixes the ordinary singularity type of the envelope: it agrees with the normalized ideal envelope up to bounded functions.

Let $\theta$ be a smooth real $(1,1)$-form on $X$, let $\varphi\in\PSH(X,\theta)$, let $K\subseteq X$ be compact and non-pluripolar, and let $w\in C(K,\mathbb R)$. Use the multiplier-ideal order characterized in \cref{thm:local-i-inputs} and the quasi-everywhere convention of \cref{sec:obstacle-envelopes}. Define
\begin{equation}\label{eq:hermitian-partial-i-envelope}
 \begin{aligned}
 E_{\theta,K}[\varphi]_{\mathcal I}(w)
 &=\Bigl(\sup\bigl\{a\in\PSH(X,\theta):
       a\preceq_{\mathcal I}\varphi,\\
 &\hspace{52pt}a\leq w\text{ quasi-everywhere on }K\bigr\}\Bigr)^*.
 \end{aligned}
\end{equation}

\begin{theorem}\label{thm:hermitian-partial-i-envelope}
The function $E=E_{\theta,K}[\varphi]_{\mathcal I}(w)$ is a genuine $\theta$-plurisubharmonic function and the largest member of its defining family. Let $q=P_\theta[\varphi]_{\mathcal I}$, and let $C$ be any upper bound from \cref{lem:compact-obstacle-upper-bound}. Then
\[
 q+\min_Kw\leq E\leq q+C.
\]
In particular, $E$ has the ordinary singularity type of $q$, and
\[
 E\sim_{\mathcal I}\varphi.
\]
\end{theorem}
\begin{proof}
\Cref{thm:hermitian-i-projection} gives $q\leq0$ and $q\sim_{\mathcal I}\varphi$. Thus $q+\min_Kw$ belongs to the defining family. \Cref{lem:compact-obstacle-upper-bound} bounds the entire family above by $C$, so its regularized supremum is genuine and gives the lower inequality.

For each competitor $a$, the upper bound $\sup_Xa\leq C$ and \cref{prop:p-below-i,prop:i-envelope-operations} give
\[
 a-C\leq a-\sup_Xa\leq P_\theta[a]_{\mathcal I}
                    \leq P_\theta[\varphi]_{\mathcal I}=q.
\]
Taking regularized suprema proves $E\leq q+C$.

The two bounds show that $E$ and $q$ differ by a bounded amount. They therefore have the same multiplier ideals, and these equal those of $\varphi$.

It remains to check the obstacle after regularization. Choose a countable subfamily by Choquet's lemma, and replace it by successive finite maxima. Finite maxima remain admissible, by \cref{thm:local-i-inputs} and by taking the finite union of the obstacle's exceptional sets. Let $a_j$ be the resulting increasing sequence. Outside a countable union of pluripolar subsets of $K$, every $a_j$ is at most $w$. The unregularized limit is therefore at most $w$ there. The regularized limit equals $E$ and changes this limit only on a further pluripolar set, by the local supremum theorem in \cref{sec:obstacle-envelopes}. Thus $E\leq w$ quasi-everywhere on $K$.

We have proved that $E$ is a member of its defining family. It dominates each competitor by construction, so it is the largest member; in particular the unregularized supremum in \eqref{eq:hermitian-partial-i-envelope} already equals $E$ at every point.
\end{proof}

\Needspace{25\baselineskip}
\begin{proposition}\label{prop:partial-envelope-calculus}
The envelope depends only on the multiplier-ideal class of $\varphi$. It is increasing in the obstacle and in the multiplier-ideal order of $\varphi$. For $w,w'\in C(K)$ and $c\in\mathbb R$,
\[
 \begin{split}
 E_{\theta,K}[\varphi]_{\mathcal I}(w+c)
       &=E_{\theta,K}[\varphi]_{\mathcal I}(w)+c,\\
 \bigl|E_{\theta,K}[\varphi]_{\mathcal I}(w)
       -E_{\theta,K}[\varphi]_{\mathcal I}(w')\bigr|
       &\leq\|w-w'\|_{C(K)}.
 \end{split}
\]
The second line means the two corresponding pointwise inequalities, including at poles. For $0<t<1$, smooth real $(1,1)$-backgrounds $\theta_i$ on $X$, for $i=0,1$, potentials $\varphi_i\in\PSH(X,\theta_i)$ and obstacles $w_i\in C(K)$, put $\theta_t=(1-t)\theta_0+t\theta_1$, and similarly define $\varphi_t,w_t$. Then
\[
 (1-t)E_{\theta_0,K}[\varphi_0]_{\mathcal I}(w_0)
       +tE_{\theta_1,K}[\varphi_1]_{\mathcal I}(w_1)
 \leq E_{\theta_t,K}[\varphi_t]_{\mathcal I}(w_t).
\]
For a smooth real $f$, changing the background and obstacle gives the exact identity
\[
 E_{\theta+\ddc f,K}[\varphi-f]_{\mathcal I}(w-f|_K)
       =E_{\theta,K}[\varphi]_{\mathcal I}(w)-f.
\]
\end{proposition}
\begin{proof}
Equality or inclusion of the defining families proves the first two assertions. Adding a constant gives a bijection of competitor families, proving translation. Apply monotonicity to $w'\pm\|w-w'\|_{C(K)}$ to obtain the two inequalities in the second line.

For concavity, the left side is $\theta_t$-plurisubharmonic and obeys the obstacle $w_t$ outside the union of two pluripolar exceptional sets. Each of its summands has the multiplier-ideal class of its prescribed singularity by the theorem. Addition and positive scaling in \cref{lem:hermitian-i-calculus} give the multiplier-ideal class of $\varphi_t$. Thus the left side is an admissible competitor on the right. Finally, $a\mapsto a-f$ is a bijection of the two families in the gauge identity, and regularization commutes with subtraction of a smooth function.
\end{proof}

The envelope bounds candidate limits of the partial Fubini--Study potentials. The next argument proves this upper bound; identifying a kernel measure with its Monge--Amp\`ere measure will require the geometric convergence hypotheses developed afterward.
